\section{Funkcja główna}
Funkcja \texttt{main} stanowi punkt centralny aplikacji, spinający 
w~jedną całość wszystkie opisane wcześniej moduły w~ramach 
głównej pętli programu, zorientowanej na ograniczenie 
zużycia energii. Po jednorazowej inicjalizacji peryferiów 
(ADC, magistrali I2C, nadajnika radiowego OOK oraz obiektu 
\texttt{PowerManager}), a~także uzbrojeniu liczników 
sprzętowych (WDT oraz timera nadajnika radiowego), urządzenie 
przechodzi we właściwy cykl pracy, w~którym większość czasu 
spędza w~stanie uśpienia — wybudzane cyklicznie jedynie przez 
przerwanie WDT, aż \texttt{sleep::ready()} zasygnalizuje upływ 
zadanego okresu (~1 minuty).

Po każdym wybudzeniu następuje krótkotrwała aktywacja modułów 
o~podwyższonym poborze prądu (za pośrednictwem metody 
\texttt{power\_manager.activate\_power\_hungry()}) oraz, przy
braku wcześniejszej inicjalizacji, jednorazowa inicjalizacja 
czujnika klimatycznego. Odczytane dane pomiarowe są następnie 
wielokrotnie (\texttt{SEND\_DATA\_N\_TIMES}) nadawane przez transmiter radiowy, 
przy czym błędy zarówno inicjalizacji, jak i~odczytu czujnika są 
świadomie ignorowane (\texttt{continue}/pusty blok \texttt{Err}) — co jest podejściem 
uzasadnionym w~kontekście urządzenia bateryjnego, pracującego bez 
nadzoru, gdzie priorytetem jest kontynuacja działania kosztem 
utraty pojedynczego pomiaru. Po zakończeniu transmisji zasilanie 
modułów energochłonnych jest ponownie wyłączane, a~stan diody 
statusowej sygnalizuje aktywność urządzenia w~trakcie całego 
cyklu pomiarowego.

Osobno zdefiniowana została procedura obsługi paniki (panic\_handler), 
która — z uwagi na brak możliwości bezpiecznego kontynuowania działania
programu — sygnalizuje błąd za pomocą migającej diody LED, a~następnie 
wymusza twardy restart mikrokontrolera poprzez skonfigurowanie WDT 
w~trybie resetu sprzętowego z~minimalnym możliwym czasem oczekiwania.